\chapter{Fungsi dan Variasinya}\label{ch:g11:func}

Bab ini menata apa saja yang perlu diketahui tentang \omterm{def:g11:func:function}{fungsi} sebelum masuk
ke kalkulus: keluarga kecil \emph{\omterm{def:g11:func:reference}{fungsi acuan}} yang \omterm{def:g10:functions:graph}{grafiknya} wajib
dihafal, makna tepat kata \emph{\omterm{def:g11:func:monotone}{naik}} dan
\emph{turun}, serta bagaimana variasi berperilaku ketika \omterm{def:g11:func:function}{fungsi} dipadukan.

\section{Fungsi dan kurva}

\begin{definition}[Fungsi, daerah asal, kurva]\label{def:g11:func:function}
Sebuah \emph{fungsi}\index{fungsi} $f$ memasangkan setiap bilangan $x$ dari
himpunan $D \subseteq \R$ (\emph{daerah asal}\index{daerah asal}nya) dengan
tepat satu bilangan $f(x)$. \emph{Kurvanya} (atau \omterm{def:g10:functions:graph}{grafiknya}) adalah himpunan
titik $\bigl(x, f(x)\bigr)$ untuk $x \in D$.
\end{definition}

\begin{example}
Rumus $f(x) = \frac{\sqrt{x+2}}{x-1}$ menuntut $x + 2 \geq 0$ (untuk
akar kuadratnya) dan $x \neq 1$ (penyebutnya taknol): \omterm{def:g11:func:function}{daerah asalnya}
$\intco{-2}{1} \cup \intoo{1}{+\infty}$.
\end{example}

\section{Fungsi acuan}

Kelima \omterm{def:g11:func:function}{fungsi} berikut beserta kurvanya harus menjadi refleks.

\begin{definition}[Fungsi acuan]\label{def:g11:func:reference}
\index{fungsi acuan}
\[
x \mapsto x^2 \text{ (pangkat dua)}, \quad
x \mapsto x^3 \text{ (pangkat tiga)}, \quad
x \mapsto \sqrt{x}, \quad
x \mapsto \frac{1}{x}, \quad
x \mapsto \abs{x}.
\]
\omterm{def:g11:func:function}{Daerah asalnya} berturut-turut $\R$, $\R$, $\intco{0}{+\infty}$,
$\R \setminus \{0\}$ dan $\R$. Kurva $x \mapsto \frac1x$ adalah sebuah
\emph{hiperbola}\index{hiperbola}.
\end{definition}

\begin{omfigure}
\begin{tikzpicture}
\begin{axis}[omaxis,
  width=7cm, height=5.6cm,
  xmin=-2.3, xmax=2.3, ymin=-3.4, ymax=3.6,
  xtick={-1,1}, ytick={-1,1}]
  \addplot[omDef, thick, domain=-1.85:1.85] {x^2};
  \addplot[omThm, thick, domain=-1.5:1.5] {x^3};
  \node[omDef, font=\small, anchor=east] at (axis cs:-1.5,1.9) {$x^2$};
  \node[omThm, font=\small, anchor=east] at (axis cs:-1.32,-1.73) {$x^3$};
\end{axis}
\end{tikzpicture}%
\begin{tikzpicture}
\begin{axis}[omaxis,
  width=7cm, height=5.6cm,
  xmin=-3.4, xmax=4.4, ymin=-3.4, ymax=3.6,
  xtick={1}, ytick={1}]
  \addplot[omProp, thick, domain=-3.2:-0.31, samples=60] {1/x};
  \addplot[omProp, thick, domain=0.31:4.2, samples=60] {1/x};
  \addplot[omDef, thick, domain=0:4.2] {sqrt(x)};
  \addplot[omThm, thick, domain=-3.2:3.2] {abs(x)};
  \node[omDef, font=\small, anchor=north] at (axis cs:3.6,1.55)
    {$\sqrt{x}$};
  \node[omThm, font=\small, anchor=east] at (axis cs:-2.4,2.8)
    {$\abs{x}$};
  \node[omProp, font=\small, anchor=west] at (axis cs:0.75,2.4)
    {$\tfrac1x$};
\end{axis}
\end{tikzpicture}

{\small Kelima kurva acuan: \omterm{def:g11:quad:parabola}{parabola} $x^2$ dan pangkat tiga $x^3$
(kiri); $\sqrt x$, $\abs x$ dan \omterm{def:g11:func:reference}{hiperbola} $\frac1x$ (kanan).}
\end{omfigure}

\section{Variasi pada sebuah interval}

\begin{definition}[Naik, turun]\label{def:g11:func:monotone}
Misalkan $f$ terdefinisi pada sebuah \omterm{def:g10:numbers:interval}{interval} $I$. \omterm{def:g11:func:function}{Fungsi} $f$ disebut
\emph{naik}\index{fungsi naik} pada $I$ bila \omterm{def:g11:func:function}{fungsi} itu mempertahankan urutan:
\[
\text{untuk semua } u, v \in I: \quad u < v \implies f(u) \leq f(v),
\]
dan disebut \emph{turun} bila \omterm{def:g11:func:function}{fungsi} itu membalik urutan ($u < v \implies
f(u) \geq f(v)$). Dengan ketaksamaan tegas di ruas kanannya, $f$ disebut
\emph{tegas} naik (masing-masing turun). \omterm{def:g11:func:function}{Fungsi} yang naik atau
turun pada $I$ disebut \emph{monoton} pada $I$.
\end{definition}

\begin{method}[Membuktikan variasi dari definisinya]\label{met:g11:func:variation}
Ambil dua bilangan sebarang $u < v$ pada \omterm{def:g10:numbers:interval}{intervalnya}, lalu tentukan tanda
selisih $f(v) - f(u)$, biasanya dengan \emph{\omterm{def:g10:algebra:expand}{memfaktorkannya}}. Jika
tandanya tetap pada \omterm{def:g10:numbers:interval}{interval} itu, simpulkan.
\end{method}

\begin{proposition}[Variasi fungsi acuan]\label{prop:g11:func:refvariations}
\begin{enumerate}
  \item $x^2$ tegas turun pada $\intoc{-\infty}{0}$ dan tegas
        \omterm{def:g11:func:monotone}{naik} pada $\intco{0}{+\infty}$.
  \item $x^3$ tegas \omterm{def:g11:func:monotone}{naik} pada $\R$.
  \item $\sqrt x$ tegas \omterm{def:g11:func:monotone}{naik} pada $\intco{0}{+\infty}$.
  \item $\frac1x$ tegas turun pada $\intoo{-\infty}{0}$ dan pada
        $\intoo{0}{+\infty}$ (tetapi \emph{tidak} pada gabungannya).
  \item $\abs x$ tegas turun pada $\intoc{-\infty}{0}$ dan
        tegas \omterm{def:g11:func:monotone}{naik} pada $\intco{0}{+\infty}$.
\end{enumerate}
\end{proposition}

\begin{proof}
Misalkan $u < v$. Kita faktorkan $f(v) - f(u)$ pada setiap kasusnya.

\emph{1.} $v^2 - u^2 = (v-u)(v+u)$. Pada $\intco{0}{+\infty}$: $v - u > 0$
dan $v + u > 0$ (sebab $v > u \geq 0$), jadi $v^2 > u^2$. Pada
$\intoc{-\infty}{0}$: $v + u < 0$, jadi $v^2 < u^2$.

\emph{2.} $v^3 - u^3 = (v - u)(v^2 + uv + u^2)$, dan faktor keduanya
sama dengan $\left(v + \frac u2\right)^2 + \frac{3u^2}{4} > 0$ kecuali bila
$u = v = 0$. Jadi $v^3 > u^3$ setiap kali $v > u$.

\emph{3.} Untuk $0 \leq u < v$:
$\sqrt v - \sqrt u = \frac{(\sqrt v - \sqrt u)(\sqrt v + \sqrt u)}
{\sqrt v + \sqrt u} = \frac{v - u}{\sqrt v + \sqrt u} > 0$.

\emph{4.} Untuk $0 < u < v$:
$\frac1v - \frac1u = \frac{u - v}{uv} < 0$ sebab $u - v < 0$ dan $uv > 0$;
perhitungan yang sama berlaku untuk $u < v < 0$. Pada gabungannya
pernyataannya gagal: $-1 < 1$ tetapi $\frac{1}{-1} = -1 < 1 = \frac11$.

\emph{5.} Pada $\intco{0}{+\infty}$, $\abs x = x$ tegas \omterm{def:g11:func:monotone}{naik}; pada
$\intoc{-\infty}{0}$, $\abs x = -x$ tegas turun.
\end{proof}

\begin{remark}
Butir 4 adalah jebakan klasik: ``turun pada $\intoo{-\infty}{0}$ dan pada
$\intoo{0}{+\infty}$'' \emph{tidak} sama dengan ``turun pada
$\R \setminus \{0\}$''. Variasi hanya bermakna pada \omterm{def:g10:numbers:interval}{interval}.
\end{remark}

\begin{example}[Membandingkan tanpa menghitung]
Karena $\sqrt x$ tegas \omterm{def:g11:func:monotone}{naik}, $\sqrt{7} < \sqrt{8}$; karena
$x^2$ tegas turun pada $\intoc{-\infty}{0}$, berlaku
$(-3)^2 > (-2)^2$. Kemonotonan mengubah perbandingan antarpeta menjadi
perbandingan antarmasukan.
\end{example}

\section{Fungsi baru dari fungsi lama}

\begin{proposition}[Transformasi]\label{prop:g11:func:transforms}
Misalkan $u$ sebuah \omterm{def:g11:func:function}{fungsi} pada \omterm{def:g10:numbers:interval}{interval} $I$ dan $k \in \R$.
\begin{enumerate}
  \item $u + k$ mempunyai variasi yang \emph{sama} dengan $u$ pada $I$;
        kurvanya adalah kurva $u$ yang tergeser tegak sebesar $k$.
  \item Untuk $\lambda > 0$, $\lambda u$ mempunyai variasi yang sama dengan
        $u$; untuk $\lambda < 0$, variasinya terbalik.
  \item Jika $u \geq 0$ pada $I$, maka $\sqrt u$ mempunyai variasi yang sama
        dengan $u$.
  \item Jika $u \neq 0$ dan $u$ bertanda tetap pada $I$, maka $\frac1u$
        mempunyai variasi yang \emph{berlawanan} dengan variasi $u$.
\end{enumerate}
\end{proposition}

\begin{proof}
Misalkan $u < v$ di $I$ (dengan menyalahgunakan notasi, kita membandingkan
nilai $x$ dengan $u < v$; tulislah $s = u(x_1)$, $t = u(x_2)$ untuk \omterm{def:g10:functions:function}{peta}
$x_1 < x_2$).

\emph{1.} $(u+k)(x_2) - (u+k)(x_1) = t - s$: selisih \omterm{def:g10:functions:function}{petanya} tidak
berubah, sehingga tandanya pun tidak berubah.

\emph{2.} $\lambda t - \lambda s = \lambda(t - s)$: tandanya tetap bila
$\lambda > 0$, dan terbalik bila $\lambda < 0$.

\emph{3.} Jika $u$ \omterm{def:g11:func:monotone}{naik} dan $x_1 < x_2$, maka $0 \leq s \leq t$, dan
karena $\sqrt{}$ bersifat \omterm{def:g11:func:monotone}{naik} (\cref{prop:g11:func:refvariations}) diperoleh
$\sqrt s \leq \sqrt t$. Alasan yang sama berlaku untuk kasus turun.

\emph{4.} Jika $u$ \omterm{def:g11:func:monotone}{naik}, $0 < s \leq t$ (atau $s \leq t < 0$), maka
$\frac1s \geq \frac1t$ sebab $x \mapsto \frac1x$ turun pada masing-masing
sisi $0$: jadi urutannya terbalik.
\end{proof}

\begin{example}\label{ex:g11:func:composite}
Telaahlah $f(x) = \sqrt{x^2 + 1}$ pada $\R$. \omterm{def:g11:func:function}{Fungsi} dalamnya
$u(x) = x^2 + 1$ bernilai positif, turun pada $\intoc{-\infty}{0}$ dan
\omterm{def:g11:func:monotone}{naik} pada $\intco{0}{+\infty}$ (kuadrat yang tergeser). Menurut butir 3, $f$
mempunyai variasi yang sama: turun lalu \omterm{def:g11:func:monotone}{naik}, dengan \omterm{def:g10:functions:extrema}{minimum}
$f(0) = 1$.
\end{example}

\begin{omfigure}
\begin{tikzpicture}
\begin{axis}[omaxis,
  width=8.6cm, height=5.8cm,
  xmin=-2.6, xmax=2.6, ymin=-1.7, ymax=4.6,
  xtick={-1,1}, ytick={1,2}]
  \addplot[omDef, thick, domain=-1.9:1.9] {x^2};
  \addplot[omThm, thick, domain=-1.55:1.55] {x^2 + 1};
  \addplot[omProp, thick, domain=-2.4:2.4, samples=80]
    {sqrt(x^2 + 1)};
  \node[omDef, font=\small, anchor=west] at (axis cs:1.9,3.4) {$x^2$};
  \node[omThm, font=\small, anchor=west] at (axis cs:1.35,3.3)
    {$x^2{+}1$};
  \node[omProp, font=\small, anchor=north] at (axis cs:-1.9,1.95)
    {$\sqrt{x^2+1}$};
\end{axis}
\end{tikzpicture}

{\small Transformasi sedang bekerja: menggeser $x^2$ ke atas sebesar $1$
(merah) mempertahankan variasinya; menarik akar kuadratnya (jingga)
mempertahankannya lagi, seperti pada \cref{ex:g11:func:composite}.}
\end{omfigure}

\section{Fungsi genap dan fungsi ganjil}

\begin{definition}[Paritas]\label{def:g11:func:parity}
Misalkan $f$ terdefinisi pada \omterm{def:g11:func:function}{daerah asal} $D$ yang setangkup terhadap $0$
(artinya $-x \in D$ setiap kali $x \in D$). \omterm{def:g11:func:function}{Fungsi} $f$ disebut
\emph{genap}\index{fungsi genap} bila $f(-x) = f(x)$ untuk semua $x \in D$,
dan \emph{ganjil}\index{fungsi ganjil} bila $f(-x) = -f(x)$ untuk semua $x \in D$.
\end{definition}

\begin{proposition}[Makna geometrisnya]\label{prop:g11:func:paritysym}
Kurva \omterm{def:g11:func:parity}{fungsi genap} simetris terhadap sumbu $y$; kurva
\omterm{def:g11:func:parity}{fungsi ganjil} simetris terhadap titik asal.
\end{proposition}

\begin{proof}
Jika $f$ \omterm{def:g11:func:parity}{genap} dan $(x, y)$ terletak pada kurvanya (jadi $y = f(x)$), maka
$f(-x) = f(x) = y$: titik cerminnya $(-x, y)$ pun terletak pada kurvanya. Jika
$f$ \omterm{def:g11:func:parity}{ganjil}, $f(-x) = -y$: titik $(-x, -y)$, cerminan $(x,y)$ terhadap
titik asal, terletak pada kurvanya.
\end{proof}

\begin{example}
$x^2$ dan $\abs x$ \omterm{def:g11:func:parity}{genap}; $x^3$ dan $\frac1x$ \omterm{def:g11:func:parity}{ganjil};
$f(x) = x^2 + x$ bukan keduanya: $f(-1) = 0$ padahal $f(1) = 2$, jadi
$f(-1) \neq \pm f(1)$. Paritas memangkas separuh pekerjaan: \omterm{def:g11:func:parity}{fungsi genap} atau
\omterm{def:g11:func:parity}{ganjil} hanya perlu ditelaah pada $\intco{0}{+\infty}$.
\end{example}

\section{Latihan}

\begin{exercise}[$\star$]\label{exo:g11:func:1}
Berikan \omterm{def:g11:func:function}{daerah asal} setiap \omterm{def:g11:func:function}{fungsi} berikut:
\[
f(x) = \sqrt{3 - x}, \qquad
g(x) = \frac{1}{x^2 - 4}, \qquad
h(x) = \frac{\sqrt{x}}{x - 5} .
\]
\end{exercise}

\begin{exercise}[$\star$]\label{exo:g11:func:2}
Dengan memakai variasi \omterm{def:g11:func:reference}{fungsi acuan}, bandingkan tanpa
kalkulator:
\[
\sqrt{11} \text{ dan } \sqrt{13}; \qquad
(-2.1)^2 \text{ dan } (-2.05)^2; \qquad
\frac{1}{0.9} \text{ dan } \frac{1}{1.1}; \qquad
(-1.01)^3 \text{ dan } (-0.99)^3 .
\]
\end{exercise}

\begin{exercise}[$\star$]\label{exo:g11:func:3}
Tentukan apakah setiap \omterm{def:g11:func:function}{fungsi} berikut \omterm{def:g11:func:parity}{genap}, \omterm{def:g11:func:parity}{ganjil}, atau bukan keduanya:
\[
f(x) = 3x^4 - x^2, \qquad
g(x) = x^3 + x, \qquad
h(x) = x^2 + x^3, \qquad
k(x) = \frac{x}{x^2+1} .
\]
\end{exercise}

\begin{exercise}[$\star$]\label{exo:g11:func:4}
Sebuah \omterm{def:g11:func:function}{fungsi} $f$ mempunyai perilaku variasi berikut: \omterm{def:g11:func:monotone}{naik} pada
$\intcc{-3}{1}$ dari $f(-3) = -2$ ke $f(1) = 4$, lalu turun pada
$\intcc{1}{5}$ sampai $f(5) = 0$. Berapa penyelesaian yang dipunyai \omterm{def:g10:algebra:equation}{persamaan}
$f(x) = 1$ pada $\intcc{-3}{5}$? Lalu $f(x) = 5$?
\end{exercise}

\begin{exercise}[$\star\star$]\label{exo:g11:func:5}
Dengan \cref{met:g11:func:variation} (tanda $f(v) - f(u)$), buktikan bahwa
$f(x) = x^2 - 4x$ tegas \omterm{def:g11:func:monotone}{naik} pada $\intco{2}{+\infty}$.
\end{exercise}

\begin{exercise}[$\star\star$]\label{exo:g11:func:6}
Tanpa menghitung turunan apa pun, berikan variasi
\[
f(x) = \sqrt{5 - x} \ \text{ di } \intoc{-\infty}{5},
\quad
g(x) = \frac{1}{x^2 + 1} \ \text{ di } \intco{0}{+\infty},
\quad
h(x) = 3 - 2\sqrt{x} \ \text{ di } \intco{0}{+\infty}.
\]
\end{exercise}

\begin{exercise}[$\star\star$]\label{exo:g11:func:7}
Misalkan $f(x) = (x-2)^2 - 3$ pada $\R$.
\begin{enumerate}
  \item Berikan variasi $f$ dan \omterm{def:g10:functions:extrema}{minimumnya}.
  \item Simpulkan variasi $g(x) = \frac{1}{(x-2)^2+1}$ dan
        \omterm{def:g10:functions:extrema}{maksimumnya}. (Perhatikan $g = \frac{1}{f + 4}$.)
\end{enumerate}
\end{exercise}

\begin{exercise}[$\star\star$]\label{exo:g11:func:8}
Kurva sebuah \omterm{def:g11:func:function}{fungsi} $f$ melalui titik $(0, 1)$, $(2, 3)$
dan $(4, 1)$, serta $f$ \omterm{def:g11:func:monotone}{naik} pada $\intcc{0}{2}$ dan turun pada
$\intcc{2}{4}$. Sketsalah sebuah kurva yang mungkin, lalu sketsalah kurva
$f + 2$, kurva $-f$, dan kurva $2f$.
\end{exercise}

\begin{exercise}[$\star\star$]\label{exo:g11:func:9}
Tunjukkan bahwa \omterm{def:g11:func:function}{fungsi} $f(x) = \frac{2x+1}{x+1}$ dapat ditulis
$f(x) = 2 - \frac{1}{x+1}$, lalu simpulkan variasinya pada
$\intoo{-1}{+\infty}$.
\end{exercise}

\begin{exercise}[$\star\star$]\label{exo:g11:func:10}
Benar atau salah (berikan alasan atau contoh penyangkal):
\begin{enumerate}
  \item Jumlah dua \omterm{def:g11:func:monotone}{fungsi naik} pada $I$ bersifat \omterm{def:g11:func:monotone}{naik} pada $I$.
  \item Hasil kali dua \omterm{def:g11:func:monotone}{fungsi naik} pada $I$ bersifat \omterm{def:g11:func:monotone}{naik} pada
        $I$.
  \item Jika $f$ \omterm{def:g11:func:monotone}{naik} pada $I$, maka $f^2$ ($= f \times f$) \omterm{def:g11:func:monotone}{naik}
        pada $I$.
\end{enumerate}
\end{exercise}

\begin{exercise}[$\star\star\star$]\label{exo:g11:func:11}
Misalkan $f(x) = x + \frac{1}{x}$ untuk $x > 0$.
\begin{enumerate}
  \item Tunjukkan bahwa untuk $0 < u < v$ berlaku
        $f(v) - f(u) = (v - u)\,\frac{uv - 1}{uv}$.
  \item Simpulkan bahwa $f$ tegas turun pada $\intoc{0}{1}$ dan
        tegas \omterm{def:g11:func:monotone}{naik} pada $\intco{1}{+\infty}$.
  \item Simpulkan bahwa $x + \frac1x \geq 2$ untuk semua $x > 0$, dengan
        kesamaan hanya di $x = 1$.
\end{enumerate}
\end{exercise}

\section{Soal: Pabrik kesimetrian}

\begin{problem}[{Soal akhir pekan --- setiap fungsi adalah bagian
genap ditambah bagian ganjil: aturan paritas, permainan transformasi, dan
minimum yang direbut tanpa kalkulus}]
\label{pb:g11:func:1}
Ada \omterm{def:g11:func:function}{fungsi} yang setangkup cermin, ada yang setangkup setengah
putaran, dan kebanyakan bukan keduanya --- namun sebuah keajaiban aljabar
kecil mengatakan bahwa \emph{setiap} \omterm{def:g11:func:function}{fungsi}, dengan tepat satu cara, adalah
jumlah sepotong bagian setangkup cermin dan sepotong bagian setangkup
setengah putaran. Soal ini membuktikan keajaiban itu, memainkan permainan
transformasi pada kurva acuan (\cref{prop:g11:func:transforms}), lalu
ditutup dengan teknik pemfaktoran selisih pada
\cref{exo:g11:func:11}, yang memenangkan soal peminimuman satu bab penuh
sebelum turunan tiba.

\medskip
\noindent\textbf{Bagian I --- Kefasihan dengan paritas.}
\begin{enumerate}
  \item \omterm{def:g11:func:parity}{Genap}, \omterm{def:g11:func:parity}{ganjil}, atau bukan keduanya (\cref{def:g11:func:parity})?
        Berikan alasan masing-masing: $x^4 - 3x^2$; \ $x^3 - x$; \ $x^2 + x$; \
        $\abs x$; \ $\frac1x$.
  \item Nyatakan makna geometris setiap putusannya
        (\cref{prop:g11:func:paritysym}): kesimetrian apa yang dipunyai
        kurva \omterm{def:g11:func:parity}{fungsi genap}? Kurva \omterm{def:g11:func:parity}{fungsi ganjil}?
  \item Paritas berkalian seperti tanda: buktikan bahwa hasil kali
        dua \omterm{def:g11:func:parity}{fungsi ganjil} bersifat \omterm{def:g11:func:parity}{genap}, lalu lengkapkan seluruh
        tabel perkaliannya (\omterm{def:g11:func:parity}{genap} $\times$ \omterm{def:g11:func:parity}{genap}, \omterm{def:g11:func:parity}{genap} $\times$
        \omterm{def:g11:func:parity}{ganjil}, \omterm{def:g11:func:parity}{ganjil} $\times$ \omterm{def:g11:func:parity}{ganjil}).
  \item Akibat singkatnya: jika $f$ \omterm{def:g11:func:parity}{ganjil} dan terdefinisi di $0$,
        berapa $f(0)$? Jika $f$ \omterm{def:g11:func:parity}{genap} dengan $f(3) = 7$, berapa
        $f(-3)$? Lalu tunjukkan bahwa satu-satunya \omterm{def:g11:func:function}{fungsi} yang sekaligus \omterm{def:g11:func:parity}{genap}
        dan \omterm{def:g11:func:parity}{ganjil} adalah \omterm{def:g11:func:function}{fungsi} nol.
  \item Berapa bagian kurva \omterm{def:g11:func:parity}{fungsi genap} yang sudah menentukan
        seluruhnya? Berikan sebuah \omterm{def:g11:func:function}{fungsi} yang bukan \omterm{def:g11:func:parity}{genap} maupun \omterm{def:g11:func:parity}{ganjil}, lalu
        jelaskan apa yang gagal.
\end{enumerate}

\medskip
\noindent\textbf{Bagian II --- Teorema penguraian.}
Untuk sebuah \omterm{def:g11:func:function}{fungsi} $f$ pada \omterm{def:g11:func:function}{daerah asal} yang setangkup terhadap $0$, tulislah
\[
E(x) = \frac{f(x) + f(-x)}{2},
\qquad
O(x) = \frac{f(x) - f(-x)}{2} .
\]
\begin{enumerate}[resume]
  \item Hitunglah $E$ dan $O$ untuk $f(x) = x^3 + x^2 + 1$, lalu
        periksa: $E$ \omterm{def:g11:func:parity}{genap}, $O$ \omterm{def:g11:func:parity}{ganjil}, $E + O = f$.
  \item Buktikan pernyataan umumnya: untuk sebarang $f$, \omterm{def:g11:func:function}{fungsi}
        $E$ \omterm{def:g11:func:parity}{genap}, $O$ \omterm{def:g11:func:parity}{ganjil}, dan $f = E + O$ --- \emph{setiap
        \omterm{def:g11:func:function}{fungsi} terbelah menjadi bagian \omterm{def:g11:func:parity}{genap} dan bagian \omterm{def:g11:func:parity}{ganjil}}.
  \item Buktikan bahwa pembelahan itu tunggal: jika
        $f = E_1 + O_1 = E_2 + O_2$ dengan $E_i$ \omterm{def:g11:func:parity}{genap} dan $O_i$
        \omterm{def:g11:func:parity}{ganjil}, pakailah pertanyaan 4 pada $E_1 - E_2$.
  \item Uraikan $f(x) = \dfrac{1}{1 - x}$ (untuk
        $x \neq \pm 1$): letakkan $E$ dan $O$ pada penyebut
        bersama $1 - x^2$ lalu sederhanakan.
  \item Sekilas ke depan, masing-masing satu kalimat: bagian \omterm{def:g11:func:parity}{genap} dan \omterm{def:g11:func:parity}{ganjil}
        \omterm{def:g11:func:function}{fungsi} eksponen kelak menjadi sepasang \omterm{def:g11:func:function}{fungsi} kembar
        tahun berikutnya (kosinus dan sinus hiperbolik); dan
        naluri yang sama --- membelah sebuah objek menjadi potongan yang
        setangkup lalu memperlakukan masing-masing dengan kesimetriannya
        sendiri --- menggerakkan analisis bunyi dan isyarat pada jilid
        universitas. Mengapa pembelahan semacam itu \emph{berguna}, secara umum?
\end{enumerate}

\medskip
\noindent\textbf{Bagian III --- Permainan transformasi.}
\begin{enumerate}[resume]
  \item Berangkat dari kurva $x \mapsto x^2$, perikan
        secara tepat kurva $(x - 3)^2$, \ $x^2 + 2$, \
        $(x - 3)^2 + 2$, \ $-x^2$, \ $2x^2$
        (\cref{prop:g11:func:transforms}).
  \item Susutkan $f(x) = \abs{x - 2} + \abs{x + 2}$ menjadi
        rumus per selang pada ketiga \omterm{def:g10:numbers:interval}{interval} yang dipotong oleh $-2$
        dan $2$, lalu perikan kurvanya (kurva itu punya dasar rata yang
        termasyhur).
  \item Transformasi mana yang menghormati kegenapan? Putuskan, dengan
        bukti atau contoh penyangkal: pergeseran mendatar
        $f(x - 3)$; pergeseran tegak $f(x) + 2$; peregangan tegak
        $2 f(x)$ (semuanya diterapkan pada $f$ yang \omterm{def:g11:func:parity}{genap}).
  \item Selesaikan $\abs{x - 2} + \abs{x + 2} = 6$ dengan
        rumus per selang pada pertanyaan 12.
  \item Merancang balik: sebuah \omterm{def:g11:quad:parabola}{parabola} bertitik puncak $(2, -3)$
        dan melalui $(0, 1)$. Susunlah kembali \omterm{def:g10:algebra:equation}{persamaannya} dalam
        \omterm{thm:g11:quad:canonical}{bentuk kanonik}.
\end{enumerate}

\medskip
\noindent\textbf{Bagian IV --- \omterm{def:g10:functions:extrema}{Minimum} tanpa kalkulus.}
\begin{enumerate}[resume]
  \item Cerminkan \cref{exo:g11:func:11} untuk
        $g(x) = x + \dfrac4x$ pada $\intoo{0}{+\infty}$: tunjukkan
        $g(v) - g(u) = (v - u)\,\dfrac{uv - 4}{uv}$, tentukan letak
        \omterm{def:g10:functions:extrema}{minimumnya}, lalu berikan nilainya.
  \item Terapkan: di antara semua persegi panjang berluas $16$, mana yang
        kelilingnya paling kecil? (Tulislah $P = 2\left(x +
        \frac{16}{x}\right)$ lalu pakai mesin pertanyaan 16.) Seorang
        kawan lama menutup jawabannya.
  \item Tentukan variasi
        $h(x) = \sqrt{x^2 + 1}$ pada seluruh $\R$, dengan memadukan alasan
        komposisi pada $\intco{0}{+\infty}$ dan alasan
        paritas untuk sisanya.
  \item Titik terdekat: carilah titik pada kurva
        $y = \sqrt x$ yang terdekat ke $(3, 0)$. (Minimumkan
        \emph{kuadrat} jaraknya --- sebuah bentuk kuadrat dalam $x$ --- lalu
        jelaskan mengapa meminimumkan kuadratnya sah.)
  \item Penutup, keliling pabriknya: kurva acuan adalah bahan
        mentahnya; transformasi adalah mesinnya; paritas dan
        penguraian adalah kendali mutunya; pemfaktoran selisih
        adalah bangku ukurnya. Satu kalimat untuk masing-masing --- lalu
        sebutkan perkakas ampuh yang diantarkan bab berikutnya.

\end{enumerate}
\end{problem}
